% Related Work -- draft v1 (2026-08-24)
% \cite keys are placeholders; fill in refs.bib. Key -> paper mapping at bottom.

\section{Related Work}
\label{sec:related}

\subsection{Explainability of Medical LLMs}
\label{sec:related-explainability}

As large language models enter clinical decision support, accuracy alone is
not the bar they must clear. Beyond being correct, a deployed model must be
inspectable: hallucination, biased outputs, opaque provenance, and
performance drift are recognized safety concerns
\cite{wang2024safety, asgari2025hallucination, ethics2025review, chen2023drift},
and both clinicians and regulators increasingly require an account of
\emph{why} a model answered as it did before its answers can be trusted
\cite{clinicalcode2025, clinicalreasoningfails2026}. For diagnosis the
requirement is concrete, because clinically real context routinely moves
answers: a referring physician's suspicion, a colleague's remark, a
patient's worry. In human diagnosticians this failure has a name ---
anchoring, a recognized contributor to diagnostic error
\cite{croskerry2003} --- and it enters LLM workflows through the same
channels \cite{mahajan2025bias}.

Explainability in medical AI has so far meant two families of methods. For
predictive models, post-hoc input attribution dominates --- SHAP, LIME, and
saliency maps assign importance to input features
\cite{lundberg2017shap, ribeiro2016lime, xaihealthcare2026}, with concept
bottlenecks providing clinician-vocabulary intermediates in imaging
\cite{cbm2026}. These methods presuppose a scalar output and many perturbed
reruns, and do not transfer cleanly to free-text diagnosis. For generative
LLMs, the de facto explanation is instead the model's own chain of thought:
a fluent, clinician-readable self-narration, prominent enough that auditing
the reasoning trace has been proposed as the safety mechanism for clinical
deployment \cite{mahajan2025bias}.

Self-narration, however, is unfaithful precisely where an audit needs it.
In the general domain, features that demonstrably move a model's answer go
unmentioned in its reasoning \cite{turpin2023say, lanham2023measuring}, and
disclosure remains rare in reasoning-tuned models \cite{chen2025reasoning}.
The medical evidence is now direct: injected clinical cognitive biases
degrade diagnostic accuracy while the accompanying explanations do not
disclose them \cite{schmidgall2024biasmedqa, reasoningnoprotect2025}; under
causal ablation and hint injection, chain-of-thought steps do not causally
drive the predictions of closed-source medical assistants, and injected
suggestions are absorbed without acknowledgment \cite{faithfulplausible2025};
structured re-grading finds individual reasoning steps decorative ---
removable without changing the answer \cite{clinicalgraphs2026}.

Both available forms of explanation therefore fail at the question a
diagnostic audit must answer --- \emph{what caused this answer} --- the
attribution family by construction, the self-narration family empirically.
We rebuild the clinically real intervention behind this evidence in a
causally controlled form --- a placebo-controlled four-arm design whose
evidence representations are bit-identical across arms, yielding a per-case
ground truth for which answers the note actually moved --- confirm on it
that the reasoning trace cannot support the proposed audit (attribution at
chance; Section~\ref{sec:results-intervention}), and then go inside the
model.

\subsection{Reading LLM Internals: From Probes to Natural-Language Readouts}
\label{sec:related-internals}

Instruments for reading model internals are mature. Linear probes recover
task variables from hidden states but answer only the question they were
trained on \cite{belinkov2022probing}; logit- and tuned-lens methods decode
token distributions from intermediate layers
\cite{nostalgebraist2020logitlens, belrose2023tunedlens}; sparse
autoencoders decompose activations into dictionaries of pre-learned
features \cite{cunningham2023sae, bricken2023monosemanticity}. All deliver
flags or fixed concepts. A younger line verbalizes instead: Patchscopes,
SelfIE, and LatentQA prompt a model to describe hidden representations in
open-ended language
\cite{ghandeharioun2024patchscopes, chen2024selfie, pan2024latentqa}, and
natural-language autoencoders (NLA) train an activation-to-text verbalizer
with a reconstruction objective \cite{anthropic2026nla} --- the instrument
we build on. Verbalization invites a sharp criticism: the description may
reflect the verbalizer's own parametric knowledge rather than the target
activation \cite{li2026privileged}. Our instrument battery ---
counterfactual swap tracking, memorization and shuffle controls, cross-case
specificity --- is designed as a direct answer
(Section~\ref{sec:results-instrument}).

In the general domain, these instruments have recently produced results
that anticipate pieces of ours. Activation probes detect the influence of
an injected hint that the chain-of-thought rationalizes away
\cite{rationalization2026}; user opinions suppress a model's learned
knowledge in late layers while the knowledge itself survives
\cite{truthoverridden2026}; hidden states predict reasoning errors that
verbalized confidence does not admit, though steering and self-correction
interventions built on such signals fail \cite{yuan2026hidden}; and correct
answers recoverable from hidden states coexist with wrong chains
\cite{mehrafarin2026hidden}. These findings are non-medical, their signals
are probes, patches, or graph distances rather than readable statements,
and they end at detection --- the one attempted continuation into
correction reports failure.

Medical work has moved beyond localization. Fraile Navarro et al. use the
same Gemma-3-12B NLA checkpoint and layer-32 activations to show that
clinical content can survive a triage output-format failure and to predict
case-level flips \cite{frailenavarro2026internal}. Tayebi Arasteh et al.
likewise recover evidence grades from hidden states when the model's
verbalized grade is near chance \cite{tayebi2026evidence}. These are direct
convergence results, not gaps we claim to fill. Their tasks are acuity
formatting and evidence grading rather than diagnosis; neither causally
isolates a referral-note suggestion, verifies the natural-language readout
as an activation-dependent instrument, or tests a controlled correction
ladder. Probes and sparse autoencoders also localize clinical knowledge and
features \cite{adrprobing2025, alignmentresistant2025, jmirsae2026, ehrsae2026}.

We chain these pieces, in medicine, on the causal testbed of
Section~\ref{sec:related-explainability}: a verified natural-language
readout attributes which cases the note moved from a single deployed run
(Section~\ref{sec:results-attribution}, with the supervised-probe upper
bound reported alongside); a positional trajectory shows that anchoring is
a rift between a preserved internal state and the emitted answer, not an
overwriting of the state (Section~\ref{sec:results-mechanism}); the readout
renders that internal state as a statement a clinician can read
(Section~\ref{sec:results-mechanism}); and feeding the statement back
recovers the moved cases where re-showing the evidence does not,
net-positive once paired with a precise selector
(Section~\ref{sec:results-correction}). We separately replicate the
behavioural anchoring effect on real case reports whose open diagnosis
vocabulary does not admit the same fixed-class probe. Internal-state and
correction claims on those case reports require separate experiments rather
than following from that behavioural replication.

% ---------------------------------------------------------------------------
% \cite key -> paper mapping (fill refs.bib from this list; delete when done)
% ---------------------------------------------------------------------------
% wang2024safety            Safety challenges of AI in medicine in the era of LLMs. arXiv:2409.18968
% asgari2025hallucination   Framework to assess clinical safety and hallucination rates (med summarisation). npj Digit Med, s41746-025-01670-7
% ethics2025review          Systematic review of ethical considerations of LLMs in healthcare. PMC12460403
% chen2023drift             How is ChatGPT's behavior changing over time? arXiv:2307.09009  [optional]
% clinicalcode2025          Cracking the clinical code: scoping review on mechanistic interpretability in medical report generation. ScienceDirect S2950363925000377
% clinicalreasoningfails2026 Why LLMs' Clinical Reasoning Fails. medRxiv 2026.01.26.26344845
% frailenavarro2026internal  Internal Representation, Not Clinical Knowledge. arXiv:2605.29889
% tayebi2026evidence         Internal-verbalized evidence-grade dissociation. arXiv:2606.29034
% croskerry2003             The importance of cognitive errors in diagnosis... Acad Med 2003
% mahajan2025bias           Cognitive bias in clinical LLMs. npj Digit Med 8:428 (2025)
% lundberg2017shap          A unified approach to interpreting model predictions (SHAP). NeurIPS 2017, arXiv:1705.07874
% ribeiro2016lime           "Why should I trust you?" (LIME). KDD 2016, arXiv:1602.04938
% xaihealthcare2026         XAI in healthcare use-case review. Frontiers AI 2026 (frai.2026.1749527)
% cbm2026                   Concept bottleneck models in medical imaging (e.g., arXiv:2410.15446 PMLR 2026; or radiologist-guided arXiv:2605.07785)
% turpin2023say             Language models don't always say what they think. NeurIPS 2023, arXiv:2305.04388
% lanham2023measuring       Measuring faithfulness in chain-of-thought reasoning. arXiv:2307.13702
% chen2025reasoning         Reasoning models don't always say what they think. Anthropic 2025, arXiv:2505.05410
% schmidgall2024biasmedqa   BiasMedQA: cognitive biases in medical LLMs. npj Digit Med 2024, arXiv:2402.08113
% reasoningnoprotect2025    LLM reasoning does not protect against clinical cognitive biases. medRxiv 2025 (저자 확인)
% faithfulplausible2025     Faithful or just plausible? Faithfulness of closed-source LLMs in medical reasoning. NeurIPS 2025, arXiv:2603.13988
% clinicalgraphs2026        Clinical reasoning graphs. arXiv:2606.29876
% belinkov2022probing       Probing classifiers: promises, shortcomings, advances. CL 2022
% nostalgebraist2020logitlens  Interpreting GPT: the logit lens. LessWrong 2020
% belrose2023tunedlens      Eliciting latent predictions with the tuned lens. arXiv:2303.08112
% cunningham2023sae         Sparse autoencoders find highly interpretable features. arXiv:2309.08600
% bricken2023monosemanticity Towards monosemanticity. Transformer Circuits 2023
% ghandeharioun2024patchscopes Patchscopes. ICML 2024, arXiv:2401.06102
% chen2024selfie            SelfIE. ICML 2024, arXiv:2403.10949
% pan2024latentqa           LatentQA. arXiv:2412.08686
% anthropic2026nla          Natural language autoencoders produce unsupervised explanations of LLM activations. Transformer Circuits 2026
% li2026privileged          Do activation verbalization methods convey privileged information? ICML 2026, arXiv:2509.13316
% rationalization2026       Catching rationalization in the act: activation probing. arXiv:2603.17199
% truthoverridden2026       When truth is overridden: internal origins of sycophancy. AAAI 2026, arXiv:2508.02087
% yuan2026hidden            Hidden error awareness in CoT: diagnostic, not causal. arXiv:2605.09502
% mehrafarin2026hidden      When CoT fails, the solution hides in the hidden states. arXiv:2604.23351
% adrprobing2025            Probing LLM hidden states for adverse drug reaction knowledge. PMC11844579
% alignmentresistant2025    Probing hidden states for calibrated, alignment-resistant predictions. medRxiv 2025.09.17.25336018
% jmirsae2026               SAEs for mechanistic interpretability of LLMs in medicine. JMIR AI 2026, e81134
% ehrsae2026                SAE decomposition of clinical sequence model representations. arXiv:2605.04072
